% The integrated commentary of the concise study, beginning on physical
% page 5. It is organized by the questions the texts of the Mass raise, not
% by the order of the Mass, and it sets the witnesses of the three
% interpretations beside one another where they answer the same question.
% Every claim is the expansive study's; the fuller argument, its quotations
% and its apparatus remain there.

\runninghead{Detailed Commentary}
\section{The Propers: Detailed Commentary}
\zlabel{triptych:concise:commentary:start}

The three interpretations begin from three questions the texts of this Mass
raise: who the sealed are, whom the Lesson counts exactly and then shows as
a crowd no one could count; why the Gospel of the feast is the Beatitudes,
whose promises stand in two tenses; and how the Gradual's promise can be
true, when saints have gone hungry and been put to death. Each
interpretation is carried by Fathers and Doctors who answer its question
where its own text presses hardest, and each reads the rest of the Mass in
that light. Their witnesses do not speak with one voice: they disagree on
the number, the crowd and the angel of the Lesson's chapter, on the census
the sealing is read as, on the land promised to the meek, on the tense of
the Gospel's promises, and on the rich who lacked.

\subsection{Who are the sealed, and who is the crowd?}

The difficulty is the Lesson's own plain sense. The number is exact and the
tribes are Israel's, and John says he \latin{heard} the number of the
sealed before he \latin{saw} the multitude; to take the number for the
whole Church may seem to drop \latin{centum quadraginta quatuor millia} in
order to reach the crowd. The list of tribes presses the same way, since it
is not the order of Jacob's sons.

St Bede keeps the number and reads it, giving reasons. The Gentiles are
counted among the tribes by St Paul's word, ``if ye be Christ's, then are
ye Abraham's seed'' (Gal 3:29). Twelve is multiplied by twelve and raised
to thousands, ``the cube of the number ten, by which is represented the
enduring life of the Church'' (\work{Explanatio Apocalypsis} I, on 7:4).
The order is not that of earthly generation: John begins with Juda,
``confession,'' the tribe from which the Lord sprang; he leaves out Dan,
from whom, ``as it is said, Antichrist is to be born'' (Gen 49:17); and he
ends with Benjamin, ``the son of the right hand,'' the elect's place after
death is destroyed (on 7:5--8). The names between are the virtues in which the faithful make
progress: the twelve thousand of each tribe are not twelve groups of people
but twelve kinds of holiness, which together make up one Church. Of the multitude Bede gives two accounts: the vision
``announces the glory of those who are to overcome the wickedness of the
last persecution,'' and it may also be read so that after the tribes of
Israel, to whom the Gospel was first preached, John ``desires to make
mention of the salvation of the Gentiles as well'' (on 7:9). Taken
together, the two give the number as the Church's completeness and the
multitude as its extent. And having expounded the chapter's end of the life
to come, Bede adds that the white-robed multitude ``may also be understood
of the present time, when `we are saved in hope'\,'' (Rom 8:24; on 7:17).

St Augustine nowhere expounds this chapter; what he says of its number and
its crowd he says in passing, inside other arguments. Setting out
Tichonius's rule of the times, he comes to the numbers Scripture favours:
``twelve into twelve gives one hundred and forty-four, which last number is
used in the Apocalypse to signify the whole body of the saints,'' and that
number ``has not reference to times, but to men'' (\work{De doctrina
christiana} III.35.51 --- the application stands in his account of
Tichonius's rule, which the text does not separate from his own). Writing
to Vincentius that the one Church is few beside the wicked and many in
herself, he instances this vision: ``in the Apocalypse, many thousands
`which no man can number,' from every tribe and tongue, are seen clothed in
white robes, and with palms of victory'' (\work{Letter XCIII}, \S30).
Preaching on the few and the many, he tells his hearers that the few who
are saved are ``many in the society of the Angels,'' and of the great
multitude coming with white robes and palms he says, ``This is the mass of
the saints'' (\work{Sermon LXI}, 1). Short applications all three, they
stand with Bede on the two points the Lesson presses: the number signifies
the whole body of the saints, and the uncounted crowd is the one Church.

Two witnesses read the chapter against Bede on part of it. Victorinus of
Pettau, whose is the one other early Latin commentary on the Apocalypse
among the witnesses read, takes the angel ascending from the east to be
``Elias the prophet, who is the precursor of the times of Antichrist,''
sent, as Malachias says, ``to recall the Jews to the faith of the people
that succeed them'' (Mal 4:5--6; \work{Commentary on the Apocalypse}, on
7:2); where Bede sees the Lord incarnate bearing the cross, Victorinus
sees the forerunner of the last persecution. Of the number he says nothing
at the Lesson's verse; when his commentary reaches the woman fleeing into
the wilderness, he writes that ``in the last times a hundred and forty-four
thousands of men should believe on the preaching of Elias'' (on 12:6): the
definite number counts the believers of the last times, not the whole
Church's completeness. On the crowd he stands with Bede: it shows ``the
number of the elect out of all believers, who, being cleansed by baptism in
the blood of the Lamb, have made their robes white, keeping the grace which
they have received'' (on 7:9), believers ``of the Jews and of the
nations.'' The keeping of the grace received adds a note Bede does not
sound, that the white robe is a gift to be kept. Victorinus also warns that
``we must follow the meaning of those things which are prophesied,'' not
look for order in the Apocalypse. His witness stands under a caution: the
editors of the English edition report that ``the manuscripts must have been
much corrupted,'' the mention of Genseric ``proves the subsequent
interpolation of his works,'' and what survives may be ``full of the
corrections of some pious disciple of St.~Augustine''; and his commentary
passes over the sealing itself and the tribes.

St Cyprian reads the crowd as the martyrs. Exhorting Fortunatus, he quotes
the uncounted multitude in white robes, who have washed their robes in the
blood of the Lamb, and concludes that ``the Christian martyrs cannot be
numbered'' (\work{Ad Fortunatum} \S11). That is the reading Bede refuses.
On the white robes washed in the blood of the Lamb, a verse the Lesson does
not read, Bede writes: ``He speaks not of martyrs alone. They are washed in
their own blood. But the blood of Jesus, the Son of God, cleanses the whole
Church from all sin'' (on 7:14). The feast's own history stands beside the refusal: Schuster traces it to a
collective feast of martyrs in fourth-century Syria and to the Pantheon
dedication of Boniface IV, whose Lesson was this one (\work{The
Sacramentary} V, pp.~208--210; his printed reference to chapter viii is an
error for vii).

The reception keeps the disagreement alive. Lucien Fromage, continuing Dom
Prosper Guéranger's \work{Liturgical Year}, reads the sealing as a second
census: at his first coming the God-Man had the world enrolled through
Caesar Augustus, and now \latin{l'heure a sonné d'un autre recensement, qui
doit consigner au livre de vie le résultat des opérations du salut}, the
hour has struck for another census, to record in the book of life the
outcome of the work of salvation; because many withdrew by their own fault
from the first enrolment, which comprised all men in the redemption, a
second and definitive one was needed to cut the guilty off from it
(\work{L'Année liturgique}, vol.~15, p.~94, citing St Gregory the Great's
Christmas homily on the census, \work{Hom.\ in Ev.}\ 8). Schuster draws the
opposite conclusion from the size of the crowd: ``Not only the hundred and
forty-four thousand predestined sons of Abraham, but a turba magna of every
age, sex, and condition'' enter heaven, and ``It is not then so hard to be
saved, as St John writes that he could not count the number of the elect''
(p.~210) --- adding at once the seal of the Tau that Ezechiel saw set on
the foreheads of those who sigh and mourn (Ezech 9:4), read with St Paul's
``as you are partakers of the sufferings, so shall you be also of the
consolation'' (2~Cor 1:7). Bede's hundred and forty-four thousand are
already the whole Church; Schuster's are sons of Abraham, with others
added; Victorinus's are the believers of the last times; Cyprian's crowd is
the martyrs. On this chapter the Fathers and their readers do not agree,
and the first interpretation follows Bede, with Augustine on the number and
the crowd.

Under this interpretation the rest of the Mass keeps one company in view.
The Introit keeps the feast \latin{sub honore Sanctorum omnium}, and the
Collect venerates all the saints \latin{sub una celebritate} and asks mercy
\latin{multiplicatis intercessoribus}. The Gradual calls them \latin{Sancti
eius}, his saints, and in the psalm's preceding verse, which the Mass does
not sing, ``the angel of the Lord shall encamp round about them that fear
him: and shall deliver them'' (Ps 33:8) --- that angel, set beside the
Lesson's angel who holds back the four from the earth, is this
interpretation's comparison, not the psalm's or Bede's. The Alleluia calls
\latin{omnes}, all who labour and are burdened, and the multitude that no
one could number answers to that word. On Bede's first account of the
multitude, the persecuted whom the Gospel blesses are the crowd the Lesson
shows in glory, and the Gospel describes the sealed on earth, bearing the
cross openly under the world's hatred. The Communion's three promises --- to see God, to be called the children
of God, to possess the kingdom --- are the ones the Lesson's crowd is shown
to enjoy; the assembly that hears them sung over the Sacrament is the same
Church as that multitude, still on the way. The Postcommunion's
\latin{fideles populi} are that Church on earth; the Sunday's prayers
beside them ask to be loosed from sins and kept from human dangers.

\subsection{Why the Beatitudes, and in what tense?}

Schuster writes that on this feast ``the Gospel can be no other than that
of the Beatitudes,'' because ``all are included in it and to each is
granted a special blessing'' (p.~210). The second interpretation follows
the four who read the eight sayings as one way --- but the Gospel blesses
people who lack, grieve and suffer, and its promises raise a question of
tense. The first and the eighth beatitude say \latin{ipsorum est regnum
caelorum}, ``theirs \emph{is} the kingdom''; the six between say that the
blessed \emph{shall} possess the land, be comforted, have their fill,
obtain mercy, see God, be called the children of God. That two promises
stand in the present and six in the future is this study's observation, not
the witnesses'. If every promise is postponed to the last day, the saints
the feast honours do not yet have what the Gospel promises; if every
promise is made present, the mourners are comforted already, which
experience does not show.

St Augustine reads the sayings as an ascent (page~3). The poor in spirit
are ``the humble and God-fearing,'' for ``the fear of the Lord is the
beginning of wisdom''; the soul that has submitted to God becomes meek
before the Scriptures, learns from them in what evils it is held and
mourns, labours free and so hungers for justice, receives counsel and shows
mercy, is purified in heart to contemplate the highest good, and at the
seventh step reaches wisdom, which makes peace in the whole man
(\work{De serm.\ Dom.}\ I.3). The mountain answers to Sinai: the lesser
precepts were given ``to a people who as yet required to be bound by
fear,'' the greater to a people ``now become suitable to set free by love''
(I.2). On the tense he holds that the promises ``can indeed be fulfilled
in this life, as we believe them to have been fulfilled in the case of the
apostles''; and the eighth beatitude, which returns to the first, he finds
figured in the eighth day, ``at the same time the first day'' (I.4.12).

St John Chrysostom begins in the same place and connects the sayings
differently. The poor in spirit are ``the humble and contrite in mind,''
``them who are awestruck, and tremble at the commandments of God''
(\work{Hom.\ in Matt.}\ 15.2); humility comes first because pride is ``the
root and fountain of all wickedness,'' and the Lord ``laid this law first
as a strong and safe foundation'' (15.3). The golden chain runs from humility through mourning, meekness, justice and
mercy to the one ``arrayed against dangers'' (15.9). Where Augustine's
link between the steps is an ascent of knowledge and of the Spirit's gifts
toward contemplation, Chrysostom's is a chain of dispositions, each
producing the next, and its end is readiness for persecution. They agree
that humility comes first and that persecution comes only to one already
formed by the rest. On the rewards Chrysostom mingles the tenses: the meek
are promised ``a sensible prize,'' for ``with the things spiritual'' the
Lord ``hath mingled the sensible'' (15.5); the mourners are comforted
``both here and there'' (15.4); and yet ``although He give different names
to the rewards, yet He brings all into His kingdom'' (15.7).

St Thomas Aquinas answers the tense with a distinction of the way and the
homeland (page~4). The poor in spirit have the kingdom in fact, for ``the
kingdom of God is within you'' (Lk 17:21). On the clean of heart he records
that some place the vision of God in the contemplation of this life, and
holds that the Lord promises it for the homeland. The eighth beatitude
\latin{designat perfectionem omnium praecedentium}, signifies the
perfection of all that precede it, because persecution tests them, as ``the
furnace trieth the potter's vessels'' (Ecclus 27:6); and \latin{persecutio
non facit beatum, sed eius causa}, persecution does not make a man blessed,
but its cause, \latin{propter iustitiam} --- here quoting Chrysostom
(\work{Super Matt.}\ cap.~5, lect.~1).

St Hilary of Poitiers reads the promises more strictly of the life to come.
The Lord on the mountain, \latin{in paternae scilicet majestatis positus
celsitudine}, set in the height of the Father's majesty, lays down
\latin{caelestis vitae praecepta}, the precepts of the heavenly life. The
land the meek inherit is \latin{ejus corporis, quod ipse Dominus assumpsit
habitaculum}, the body which the Lord himself took as his dwelling ---
where Chrysostom's meek receive ``a sensible prize.'' The clean of heart
see God because \latin{non enim nisi spiritu perfecti, et immortalitate
mutati \dots\ quod in Deo est immortale cernemus}: only when we are
perfected in spirit and changed by immortality shall we behold what in God
is immortal (\work{In Matt.}\ IV). The difference between Augustine and
Hilary is real: what Augustine believes fulfilled in the apostles, Hilary
reserves to the perfected. On the feast the present tense of the first and
the eighth beatitude is true of the saints in heaven in its completion, and
of the Church on earth in its beginning.

Under this interpretation the other texts of the Mass fall along the way.
The Alleluia's context in Matthew names the disposition with which the
Beatitudes begin, ``learn of me, because I am meek, and humble of heart''
(Mt 11:29), and Augustine, preaching on that verse, bids his hearers
``first think of the foundation of humility,'' in a building that reaches
``even to the Vision of God'' (\work{Sermo} 69, 1--3) --- the reward of
the sixth beatitude. The Gradual's \latin{Timete Dominum} names the fear
of the first step; the Lesson's multitude in white robes and with palms is
the way's company arrived (Bede's comment is on the Apocalypse; the join is
this interpretation's own); the Collect's merits of the saints are the way
walked, and the Secret asks honour for those who have finished it.

The Communion remains: why does the Mass, having read eight beatitudes,
sing three of them again? The witnesses do not say; Fromage's
continuation of Guéranger remarks only that the Church receives in the
Sacrament the three beatitudes she has just heard, ``the three last
beatitudes.'' The three are the ones whose promises name the way's end: in Augustine's order
the sixth and seventh steps and the eighth that returns to the first, in
Aquinas's account the vision of God and the perfection of all the rest ---
the clean of heart shall see God, the peacemakers shall be called the
children of God, and theirs, as the poor in spirit's, is the kingdom of
heaven. That the Mass sings the vision, the sonship and the kingdom where
the assembly receives the Sacrament is this interpretation's inference,
drawn from the two orders and from the Communion's place in the Mass. The
Postcommunion then asks that those who have walked the way defend those who
walk it, and the commemorated Sunday's Collect asks for release from sins
contracted \latin{pro nostra fragilitate}: the same Mass that reads the way
of the saints asks pardon for frailty.

\subsection{A promise tested by hunger and death}

The third interpretation meets the strongest objection to the Gradual where
St Augustine states it, as plainly as any objector. Many will not fear God
for fear of going hungry; and when the psalm goes on, ``The rich do lack
and suffer hunger, but they that seek the Lord shall not want any good
thing,'' then ``if thou receive it according to the letter, He seemeth to
deceive thee, for thou seest that many rich men that are wicked die in
their riches.'' The man who has seen such a rich man buried says, ``the
Scripture has deceived me,'' and adds, ``Daily I rise up to Church, daily I
bend the knee, daily I seek the Lord, and have nothing good'' (\work{En.\
in Ps.}\ 33 [NPNF ``Psalm XXXIV''], 13). Two features of the Mass bear on
the objection before the witnesses answer it. The Gradual leaves out the
words about the rich, \latin{Divites eguerunt, et esurierunt}, on which the
commentators build part of their answer; the Mass sings the promise without
the contrast. And the Offertory's verses answer the wicked who resolved to ``condemn
him to a most shameful death'' (Wis 2:20).

St Basil's answer begins from what fear does. \latin{Nisi timor instituerit
vitam nostram, sanctimonia in corpore comparari non potest}: unless fear
has ordered our life, holiness cannot be gained in the body. \latin{Timenti itaque
non est inopia}: he who is kept by fear from every unseemly act
\latin{nihil deest ad quamlibet virtutem, sed perfectus est}, lacks nothing
for any virtue, but is perfect. Of the rich who lacked, Basil offers that the psalm may call rich the
people of Israel (Rom 9:4), who came into want through their sin against
the Lord and, after they killed the bread of life (Jn 6:48), suffered a
famine ``not of bread \dots\ but of hearing the word of the Lord'' (Amos
8:11). Those from the nations
who learned to seek the Lord are not diminished in any good, and
\latin{Perfectum omnino bonum Deus ipse est}, the wholly perfect good is
God himself. Basil then puts the objection as a dilemma: either the apostles did not
seek the Lord, or Scripture is false; neither can be said, for Paul sought
the Lord and nothing good was lacking to him, though thrice beaten with
rods, once stoned, thrice shipwrecked, in hunger and thirst to the last
hour (2~Cor 11:25--27; \work{Hom.\ in Ps.}\ 33, PG 29, cols.~365--368).

St Robert Bellarmine (whose English numbers these verses one lower than the
Vulgate) reads the fear as ``the pure and filial fear that is the
characteristic of the saints and of the perfect,'' for ``the perfect lover
fears vehemently lest he may offend his beloved in any way.'' The promise
then has two senses. In ``the higher meaning,'' the essence of perfection
is to feel no want: ``what want can the friend of God, who owns
everything, feel''; ``and if the just appear sometimes to be in want, they
really are not so, because they get patience, better than any riches, to
bear it.'' In ``an humbler sense'' the
same words teach the simple ``that God will supply his servants with the
temporal things of the world, and will not desert them in time of
necessity.'' Of the next verse: those who seek the Lord ``shall not be
deprived of any good,'' ``because, as they cling to the Supreme Good, they
possess all that is good'' (\work{Explanatio}, Ps~33, his vv.~9--10). Basil
and Bellarmine agree that the promise is kept in the good the God-fearing
possess, God himself; they differ on the rich, whom Basil reads as Israel
impoverished of the word and Bellarmine, in the plain sense, as the wicked
rich reduced to poverty and, in the higher, as the covetous, who ``always
hunger and need.'' And Bellarmine keeps a temporal sense, bounded as the
humbler meaning, that Basil does not state.

St Augustine answers his own objection by naming the goods: ``When thou art
filled with spiritual riches, canst thou be poor? \dots\ art thou poor who
hast the chamber of thy heart filled with such jewelry of virtues, justice,
truth, charity, faith, endurance?'' The rich who lack, lack bread, and he names it from the
Gospel, ``I am the Living Bread which came down from Heaven,'' and with it
the beatitude, ``Blessed are they which do hunger and thirst after
righteousness: for they shall be filled'' (\S14) --- his comment on the
half-verse the Gradual does not sing, while ``what manner of good, I have
already said'' is his comment on the half it does.

St Thomas Aquinas expounds \latin{non est inopia} first of spiritual goods,
since ``the fear of the Lord is his treasure'' (Isa 33:6), then of bodily
goods: it may happen that one who fears God has little, but not that he is
needy, for \latin{Inops est qui se deficientem reputat}, he is needy who
counts himself lacking, and those who fear God are content with what they
have (Phil 4:11--12); and he adds, \latin{Item Deus quaerentibus se
subvenit in necessitate}, God also comes to the help of those who seek him
in their need. He then states the objection in Augustine's name from St
Paul, ``Even unto this hour we both hunger and thirst and are naked''
(1~Cor 4:11), and answers that \latin{Deus est nutritor et medicus}, God is
both nourisher and physician: the physician withholds food from the sick
man because it helps his health, and so God, as it helps our salvation,
\latin{quandoque inopiam immittit, quandoque divitias confert}, sometimes
sends want and sometimes gives riches (\work{Super Psalmos} 33,
nn.~10--11). The want of the God-fearing is given for their health.

Their answers differ in what they say is promised: God himself (Basil,
Bellarmine), the virtues (Augustine), and want that serves salvation
(Aquinas). Two keep a temporal sense, each in a bounded form: Bellarmine as
the humbler meaning, Aquinas as God's help in necessity. The Gradual's
omission of \latin{Divites eguerunt} does not remove the difficulty,
because the singers still know saints who were poor; it means that what the
Mass sings is the promise alone, and what the Fathers say of the rich is
said about a half-verse the Mass leaves out.

The Offertory carries the same question from want to death. Chrysostom
quotes its first clause against the wandering of the souls of the dead
(\work{Hom.\ in Matt.}\ 28), Augustine to explain ``My soul is continually
in thy hand'' (\work{En.\ in Ps.}\ 118, on v.~109), and St Gregory's just
man is left lifeless amidst the persecutors' hands (page~4). The 1962 antiphon reads \latin{tormentum malitiae}, the
torment of malice, where the Vulgate reads \latin{tormentum mortis}, the
torment of death --- its history among the Missal printings is given with
the appointed texts of the full study, and when its closing
\latin{alleluia} entered is not established --- and it leaves out the
verses in which the departure of the just ``was taken for misery.'' What
the 1962 form says is that the torment of malice will not touch the just.
Under this interpretation the Lesson shows the promise where it is kept:
the multitude before the throne, crying ``Salvation to our God,'' wants
nothing; what the psalm promises and the assembly sings, the vision shows
as fact. The Collect asks for the abundance of mercy that is desired. The
Communion names the good that is not lacking as the vision of God himself,
which Bellarmine's higher sense makes the whole of the promise.

\subsection{The Church on the way and the Church at rest}

The orations of the feast, read together, set the Church on earth beside
the Church in heaven. The Collect keeps all the saints in one celebration
and asks mercy through their multiplied intercession; the Secret asks that
the offering be pleasing for the honour of all the just and salutary to
those who offer it; the Postcommunion holds together the Introit's joy and
the Collect's defence. Schuster dwells on
its name for the Church on earth, \latin{fideles populi}, ``the people
journeying towards eternity with the eyes, and in the light of faith''
(p.~211). The feast's parts are old, though the feast in this form is not
only old: Schuster finds the Collect already in the Gelasian Sacramentary,
for a collective feast of the Apostles within the octave of Ss Peter and
Paul (p.~210), and the Introit is St Agatha's antiphon widened to all the
saints (page~3).

The three prayers of the commemorated Sunday belong to another day and are
said here because the feast has displaced the Sunday. Its Collect asks God to loose the sins of his peoples from the bonds their
frailty contracted; its Secret asks him to bring to completion what he has
given \latin{immeritis}, to the undeserving; its Postcommunion asks that
those given to rejoice in divine things be not left subject to human
dangers. Said after the feast's own
prayers, they keep the assembly in its place: the Church that venerates the
sealed is the Church that asks to be absolved, and the hands that receive
divine things still need protection from human dangers. The Preface of the
Mass is not settled. The feast has none of its own, and the general rubrics
can be read for the Preface of the Most Holy Trinity (p.~293, no.~1082) or
for the common Preface (p.~297, no.~1086); both candidates are named here
as the study names them, and neither is asserted, and no interpretation
rests on either Preface's words.

\subsection{Where the three interpretations meet and differ}

All three rest on the fear of the Lord and on humility, though they come to
them from different texts. In the first, the just rejoice ``not in
yourselves'' but in the Lord, and the seal is set on those who confess. In
the second, the poor in spirit are the humble and God-fearing, and
Augustine's way begins with the gift of fear. In the third, the Gradual's
\latin{Timete Dominum} is the subject itself: for Basil and Bellarmine the
one who fears lacks nothing, because the perfect good, God himself, is
possessed. In each interpretation holiness begins in what God
gives: the seal, the humility of the first beatitude, the fear the Gradual
commands.

They attend to different times. The first looks at the saints already
before the throne and reads back from there to the sealing on earth. The
second follows the way the saints walked, from its first step to its last.
The third stands where the assembly stands, with the promise still being
tested by hunger and death. Bede's remark that the vision of the
white-robed multitude ``may also be understood of the present time,'' and
Aquinas's distinction between the rewards begun on the way and completed in
the homeland, give the first two a bridge to the third's present.

The three meet in the same texts and read them differently. For the first,
the Lesson's multitude with palms is the Church that has overcome and the
Offertory's just are the sealed who have passed through death; for the
second, the same crowd is the way's end, its robes baptism and its palms
the cross's triumph; for the third, the Offertory answers the scandal of
the just man's death and the crowd is the Gradual's promise shown as kept.

Where the senses differ, they differ because the questions differ:
allegorically, the number is the whole Church for Bede and Augustine, a
company of the last times for Victorinus, a census of judgment for Fromage,
and the mountain is a new Sinai for Augustine and the height of the
Father's majesty for Hilary; morally, the faithful are to confess openly
and rejoice in the Lord (first), to begin in humility and be ready to
suffer for justice (second), and to keep the filial fear that lacks
nothing (third); anagogically, the elect stand at the right hand and give
unceasing praise (first), the one reward is completed in the homeland
(second), and the just are in the hand of God and in peace (third). What
the feast asks of those who hear it, the three answer together: to confess the faith openly and rejoice
in the Lord, not in themselves; to begin in humility and the fear of God
and be ready to suffer for justice; to fear the Lord with a son's fear and
to bear want patiently. The witnesses still disagree on the number, the
crowd and the angel, on the census, on the land of the meek and on the rich
who lack.
